% Generated by report/tools/tables.py from the result files of the repository. Do not edit.
\begin{table}[tbp]
\centering
\footnotesize
\setlength{\tabcolsep}{4pt}
\caption[E5 hors échantillon]{Expérience E5 vérifiée hors échantillon (écart 25 de \file{sim/SIM\_PROTOCOL.md}, écrit avant ce calcul). E5 prend pour bon ordre le signe des moyennes de l'échantillon même où il tire avec remise : sa part mesure la stabilité du signe observé, pas la probabilité de retrouver le vrai ordre. Signe de l'écart des graines 0 à 9 retrouvé sur les graines 10 à 19 : 20 comparaisons sur 24, 5 serrées sur 7 ; signes qui changent : HalfCheetah : PPO PyTorch / PPO JAX ; Hopper : CrossQ / PPO PyTorch ; Hopper : CrossQ / PPO JAX ; Swimmer : CrossQ / PPO PyTorch. Ici, $n$ essais par programme tirés sans remise dans une moitié des graines (10\,000 tirages, \texttt{default\_rng(2026)}), comptés justes quand ils retrouvent l'ordre des moyennes de l'autre moitié, dans les deux sens. Comparaison serrée : écart relatif sous 20~\% sur les graines 0 à 9. À $n = 10$, le tirage prend toute la moitié, d'où 100 ou 0~\% par comparaison : la moyenne des serrées donne alors la part des signes retrouvés. \textbf{En gras}, la moyenne des comparaisons serrées, la mesure sans circularité. Sources : \file{sim/analysis/e5\_holdout.json} (\file{sim/analysis/e5\_holdout.py}), \file{sim/analysis/e5.json} (tableau~\ref{tab:e5}).}
\label{tab:e5-hors-echantillon}
\begin{tabularx}{\linewidth}{@{}>{\raggedright\arraybackslash}Xrrr@{}}
\toprule
Mesure (part des tirages qui retrouvent l'ordre, \%) & $n = 3$ & $n = 5$ & $n = 10$ \\
\midrule
\multicolumn{4}{@{}l}{\emph{E5 préenregistré, ordre jugé sur le même échantillon}} \\
\quad Médiane, 24 comparaisons & 94 & 97 & 99 \\
\quad Médiane, 7 comparaisons serrées & 65 & 71 & 77 \\
\multicolumn{4}{@{}l}{\emph{Tirés dans les graines 0 à 9, jugés sur les graines 10 à 19}} \\
\quad Médiane, 24 comparaisons & 96 & 99 & 100 \\
\quad Médiane, 7 comparaisons serrées & 61 & 70 & 100 \\
\textbf{\quad Moyenne, 7 comparaisons serrées} & \textbf{60} & \textbf{65} & \textbf{71} \\
\multicolumn{4}{@{}l}{\emph{Tirés dans les graines 10 à 19, jugés sur les graines 0 à 9}} \\
\quad Médiane, 24 comparaisons & 92 & 99 & 100 \\
\quad Médiane, 7 comparaisons serrées & 71 & 80 & 100 \\
\textbf{\quad Moyenne, 7 comparaisons serrées} & \textbf{69} & \textbf{71} & \textbf{71} \\
\bottomrule
\end{tabularx}
\end{table}
